% =====================================================================
% Capítulo: Requisitos del Sistema
% Fuentes: propuesta del TFG (Objetivo, Descripción y Alcance),
% docs/tfg/REQUISITOS.md, DECISIONES.md (ADR-004, 005, 011, 017, 020, 021)
% y borradores/04-analisis.tex. Estado de las fuentes: F3c cerrada.
% =====================================================================

En este capítulo se presentan los objetivos de negocio y del sistema, los
actores que interactúan con PymeKit y el catálogo de requisitos funcionales,
no funcionales y de información. Al final se relaciona cada requisito con el
apartado de la propuesta del TFG del que procede y se resume la brecha entre
esos requisitos y las bases de código de partida.

Los identificadores (OBN, OBS, RF, RNF, RI) son estables: se citan en el
código mediante etiquetas \texttt{[TFG]}, en la matriz de trazabilidad del
repositorio y en esta memoria, y no se renumeran. Si un requisito cambia, se
modifica su redacción y se registra la decisión que lo justifica.

% ---------------------------------------------------------------------
\section{Objetivos de negocio}
\label{sec:req-objetivos-negocio}

Quien desarrolla un producto SaaS para pymes debe resolver, antes de escribir
su lógica de negocio, un conjunto de problemas que se repiten en casi todos
los productos de este tipo: identidad, aislamiento entre clientes, equipos y
permisos, cobro, administración interna y calidad transversal. Estos
componentes no aportan valor diferencial, pero consumen gran parte del
esfuerzo y concentran el riesgo de seguridad. De ahí se derivan los objetivos
de negocio de la tabla~\ref{tab:req-obn}.

\begin{table}[htbp]
  \centering
  \footnotesize
  \caption{Objetivos de negocio.}
  \label{tab:req-obn}
  \begin{tabularx}{\textwidth}{@{}lX@{}}
    \toprule
    \textbf{ID} & \textbf{Objetivo} \\
    \midrule
    OBN-01 & Reducir el tiempo de puesta en marcha de un nuevo producto SaaS para pymes, partiendo de una base que ya resuelve los componentes comunes. \\
    OBN-02 & Estandarizar buenas prácticas de seguridad, control de acceso y persistencia, de modo que cada producto nuevo las herede sin reimplementarlas. \\
    OBN-03 & Permitir que una pyme administre su web pública, su contenido y sus datos desde una consola, sin acceder directamente a la base de datos ni depender de un desarrollador para cada cambio. \\
    OBN-04 & Facilitar la monetización mediante suscripciones con una pasarela de pago estándar. \\
    OBN-05 & Servir como material didáctico y técnico documentado para quien reutilice la plataforma (código comentado en español, decisiones registradas y manuales). \\
    \bottomrule
  \end{tabularx}
\end{table}

Durante el desarrollo se precisó el enfoque del producto (ADR-020): PymeKit es
la base de la web de \emph{una} pyme, con una página pública, un blog
gestionado desde la consola y un inicio de sesión que puede corresponder a un
super-administrador, a personal interno o a un cliente. Los módulos SaaS
(equipos, invitaciones y suscripciones) se conservan como opcionales, porque
permiten tanto ese escenario como el de un SaaS multi-cliente.

% ---------------------------------------------------------------------
\section{Objetivos del sistema}
\label{sec:req-objetivos-sistema}

Los objetivos del sistema concretan los componentes que enumera el objetivo de
la propuesta (tabla~\ref{tab:req-obs}).

\begin{table}[htbp]
  \centering
  \footnotesize
  \caption{Objetivos del sistema y objetivos de negocio a los que contribuyen.}
  \label{tab:req-obs}
  \begin{tabularx}{\textwidth}{@{}lXl@{}}
    \toprule
    \textbf{ID} & \textbf{Objetivo del sistema} & \textbf{OBN} \\
    \midrule
    OBS-01 & Ofrecer un \textit{frontend} web con página pública, precios, blog y área privada. & 01, 03 \\
    OBS-02 & Ofrecer un \textit{backend} con API tipada y validada para la lógica de negocio y el acceso a datos. & 01 \\
    OBS-03 & Autenticar a los usuarios y autorizar cada acceso, con segundo factor y aislamiento entre cuentas aplicado en la base de datos. & 02 \\
    OBS-04 & Gestionar usuarios, cuentas de equipo, invitaciones, roles y permisos. & 01, 02 \\
    OBS-05 & Integrar una pasarela de pago para suscripciones. & 04 \\
    OBS-06 & Proporcionar una consola de administración con un CMS de datos y contenidos, con permisos y auditoría propios. & 03 \\
    OBS-07 & Ser reutilizable por configuración y estar documentado y verificado de forma automática. & 01, 05 \\
    \bottomrule
  \end{tabularx}
\end{table}

% ---------------------------------------------------------------------
\section{Actores}
\label{sec:req-actores}

\begin{description}[style=nextline]
  \item[Visitante.] Persona no identificada que consulta la web pública
        (inicio, precios y blog) y puede registrarse.
  \item[Usuario.] Persona registrada. Dispone de una cuenta personal y puede
        pertenecer a una o varias cuentas de equipo.
  \item[Propietario del equipo.] Usuario con el rol de mayor nivel en una
        cuenta de equipo; gestiona sus miembros y su suscripción.
  \item[Miembro.] Usuario de un equipo con un rol de menor nivel y, por
        tanto, con permisos limitados.
  \item[Super-administrador.] Personal del proveedor con control total de la
        plataforma y, por extensión, del CMS (ADR-014). Debe usar siempre
        segundo factor.
  \item[Personal del CMS.] Personal de la pyme o del proveedor (gestor de
        contenidos, soporte) que accede a la consola con los permisos
        limitados que le asigna el control de acceso propio del CMS
        (ADR-014, ADR-016).
  \item[Desarrollador.] Quien reutiliza PymeKit para construir un producto
        nuevo; interactúa con el repositorio, la configuración y los
        manuales, no con la interfaz.
  \item[Pasarela de pago.] Sistema externo (Stripe) que procesa los cobros y
        notifica los cambios de suscripción mediante \textit{webhooks}
        (peticiones HTTP que el sistema externo envía a la aplicación cuando
        ocurre un evento).
\end{description}

% ---------------------------------------------------------------------
\section{Requisitos funcionales}
\label{sec:req-rf}

La prioridad es \textbf{M} (imprescindible: forma parte del prototipo mínimo
que exige la propuesta) o \textbf{D} (deseable). La última columna indica el
objetivo del sistema que desarrolla cada requisito.

\begin{table}[htbp]
  \centering
  \footnotesize
  \caption{Requisitos funcionales.}
  \label{tab:req-rf}
  \begin{tabularx}{\textwidth}{@{}lXcl@{}}
    \toprule
    \textbf{ID} & \textbf{Requisito} & \textbf{Prior.} & \textbf{OBS} \\
    \midrule
    RF-01 & Interfaz web para el usuario final: página de inicio pública (\textit{landing}), precios, blog público y área privada con navegación. & M & 01 \\
    RF-02 & API de \textit{backend} para la lógica de negocio y el acceso a datos mediante funciones de servidor (\textit{server functions}) tipadas y validadas. & M & 02 \\
    RF-03 & Registro e inicio de sesión (contraseña y enlace mágico), recuperación de contraseña y verificación del correo electrónico. & M & 03 \\
    RF-04 & Segundo factor de autenticación (MFA) y proveedores OAuth configurables. & D & 03 \\
    RF-05 & Gestión de la cuenta personal: perfil, correo, contraseña y baja. & M & 04 \\
    RF-06 & Cuentas de equipo: crear equipos, invitar miembros, asignar roles y aplicar permisos por funcionalidad. & M & 04 \\
    RF-07 & Suscripciones con pasarela de pago: planes, pago (\textit{checkout}), portal de cliente y sincronización mediante \textit{webhooks}. & M & 05 \\
    RF-08 & Panel de super-administración: listar y gestionar usuarios y cuentas, bloquear, suplantar y restablecer. & M & 06 \\
    RF-09 & CMS de datos y contenidos integrado en la consola de administración: explorar y editar tablas, gestionar usuarios y almacenamiento, publicar el contenido del blog, con permisos propios. & M & 06 \\
    RF-10 & Registro de auditoría de las acciones realizadas desde el CMS. & D & 06 \\
    RF-11 & Paneles (\textit{dashboards}) configurables en el CMS. & D & 06 \\
    RF-12 & Notificaciones dentro de la aplicación y correos transaccionales. & D & 04 \\
    RF-13 & Interfaz en español, con la infraestructura de internacionalización preparada para añadir otros idiomas. & M & 07 \\
    \bottomrule
  \end{tabularx}
\end{table}

Dos requisitos han cambiado de redacción durante el proyecto, y ambos cambios
están registrados como decisiones:

\begin{itemize}
  \item \textbf{RF-09} se planteó al principio como un CMS en una aplicación
        separada que compartía la base de datos. Tras evaluar alternativas se
        decidió integrarlo en la consola de administración (ADR-011) y,
        después, usarlo también para gestionar el contenido del blog que
        publica la web (ADR-017), que es el caso de uso de «administración de
        contenidos» que pide la propuesta.
  \item \textbf{RF-13} pedía en origen una interfaz en español por defecto y
        en inglés. A petición del autor, la interfaz pasó a ofrecer
        \emph{solo} español, conservando claves, mensajes y configuración de
        idiomas como punto de extensión documentado (ADR-021).
\end{itemize}

RF-11 se identificó desde el principio como el primer candidato a recortar si
el plazo lo exigía; finalmente se implementó (ADR-019).

% ---------------------------------------------------------------------
\section{Requisitos no funcionales}
\label{sec:req-rnf}

Los requisitos no funcionales se clasifican según las características de
calidad de la norma ISO/IEC 25010~\cite{iso25010}. Cada uno se acompaña del
modo en que se verifica, para que sea comprobable y no una mera declaración
de intenciones.

\begin{table}[htbp]
  \centering
  \footnotesize
  \caption{Requisitos no funcionales, característica ISO/IEC 25010 y forma de verificarlos.}
  \label{tab:req-rnf}
  \begin{tabularx}{\textwidth}{@{}l>{\raggedright\arraybackslash}p{5.4cm}>{\raggedright\arraybackslash}p{2.2cm}X@{}}
    \toprule
    \textbf{ID} & \textbf{Requisito} & \textbf{ISO 25010} & \textbf{Verificación} \\
    \midrule
    RNF-01 & \textbf{Reutilización y configurabilidad}: la plataforma se adapta a un producto nuevo cambiando la configuración (ficheros \texttt{*.config.ts} y variables de entorno), sin tocar el núcleo. & Mantenibilidad (reusabilidad), flexibilidad & Evaluación de la reutilización: pasos y tiempo para arrancar un proyecto nuevo \\
    RNF-02 & \textbf{Seguridad}: la autorización se aplica en la base de datos (RLS) además de en la aplicación; secretos fuera del repositorio; validación de todas las entradas. & Seguridad & Pruebas pgTAP y revisiones adversariales de RLS y del código \\
    RNF-03 & \textbf{Aislamiento multi-tenant}: los datos de una cuenta no son accesibles desde otra. & Seguridad (confidencialidad) & Pruebas pgTAP de aislamiento \\
    RNF-04 & \textbf{Modularidad}: monorepo de paquetes desacoplados con dependencias explícitas y exportaciones separadas para cliente y servidor. & Mantenibilidad (modularidad) & Estructura del monorepo, comprobación de dependencias y de tipos \\
    RNF-05 & \textbf{Mantenibilidad}: TypeScript estricto, \textit{lint} y formato automáticos y código comentado en español didáctico. & Mantenibilidad (analizabilidad) & Integración continua y revisión de comentarios \\
    RNF-06 & \textbf{Calidad verificable}: pruebas unitarias, de base de datos y E2E de los casos de uso principales. & Mantenibilidad (capacidad de prueba) & Integración continua en verde \\
    RNF-07 & \textbf{Desplegabilidad}: despliegue reproducible con contenedores y en la nube, con manual de instalación. & Flexibilidad (instalabilidad) & Despliegue de prueba \\
    RNF-08 & \textbf{Usabilidad y accesibilidad}: interfaz adaptable, componentes accesibles, tema claro/oscuro y ayuda contextual en la consola. & Capacidad de interacción & Revisión manual y E2E \\
    RNF-09 & \textbf{Rendimiento razonable}: carga de datos en el servidor, caché de consultas y paginación en el CMS. & Eficiencia de desempeño & Revisión manual y pruebas puntuales \\
    \bottomrule
  \end{tabularx}
\end{table}
% VERIFICAR: nombres de las características según la edición de la norma que
% se cite (2011 o 2023; en 2023 «usabilidad» pasa a «capacidad de
% interacción» y «portabilidad» a «flexibilidad»).

% ---------------------------------------------------------------------
\section{Requisitos de información}
\label{sec:req-ri}

La tabla~\ref{tab:req-ri} recoge la información que el sistema debe almacenar.
Su modelo conceptual se desarrolla en el capítulo de Análisis
(sección~\ref{sec:ana-modelo-datos}).

\begin{table}[htbp]
  \centering
  \footnotesize
  \caption{Requisitos de información.}
  \label{tab:req-ri}
  \begin{tabularx}{\textwidth}{@{}l>{\raggedright\arraybackslash}p{3.1cm}Xl@{}}
    \toprule
    \textbf{ID} & \textbf{Información} & \textbf{Datos principales} & \textbf{RF} \\
    \midrule
    RI-01 & Usuarios & Identidad, correo, factores MFA, estado de bloqueo (gestionados por el servicio de autenticación) & 03--05 \\
    RI-02 & Cuentas & Cuenta personal o de equipo, nombre, identificador legible (\textit{slug}), propietario principal & 05, 06 \\
    RI-03 & Membresías y roles & Pertenencia de un usuario a una cuenta, rol con nivel jerárquico y permisos por funcionalidad & 06 \\
    RI-04 & Invitaciones & Correo invitado, rol ofrecido, testigo (\textit{token}), caducidad, quién invita & 06 \\
    RI-05 & Facturación & Cliente de la pasarela, suscripciones y sus líneas, pedidos de pago único & 07 \\
    RI-06 & Notificaciones & Destinatario, tipo, texto, enlace, leída o no & 12 \\
    RI-07 & Contenido del blog & Entradas (título, \textit{slug}, cuerpo en Markdown, estado, fecha de publicación, autor), categorías y etiquetas & 01, 09 \\
    RI-08 & Personal y permisos del CMS & Cuentas del CMS, roles con rango, grupos de permisos, permisos de datos, de sistema y de almacenamiento, con vigencia & 09 \\
    RI-09 & Metadatos del CMS & Por tabla: nombre visible, descripción, área de navegación, formato de visualización, configuración de columnas y diseño de la ficha; vistas guardadas & 09 \\
    RI-10 & Auditoría & Operación, tabla, registro, datos antes y después, gravedad, instantánea del autor & 10 \\
    RI-11 & Paneles & Panel, propietario, compartición por rol, \textit{widgets} y su posición en la cuadrícula & 11 \\
    \bottomrule
  \end{tabularx}
\end{table}

% ---------------------------------------------------------------------
\section{Trazabilidad con la propuesta}
\label{sec:req-trazabilidad}

La tabla~\ref{tab:req-traza-propuesta} relaciona los apartados de la
propuesta del TFG con los requisitos que los cubren y el capítulo de la
memoria donde se desarrollan. La relación inversa, de cada requisito con el
código y las pruebas que lo verifican, se presenta en el capítulo de Pruebas
(sección~\ref{sec:pru-resultados}).

\begin{table}[htbp]
  \centering
  \footnotesize
  \caption{Trazabilidad entre la propuesta del TFG, los requisitos y la memoria.}
  \label{tab:req-traza-propuesta}
  \begin{tabularx}{\textwidth}{@{}>{\raggedright\arraybackslash}X>{\raggedright\arraybackslash}p{3.6cm}>{\raggedright\arraybackslash}p{3.2cm}@{}}
    \toprule
    \textbf{Apartado de la propuesta} & \textbf{Requisitos} & \textbf{Capítulo} \\
    \midrule
    Objetivo: \textit{frontend}, \textit{backend}, autenticación, usuarios y roles, pagos y CMS & RF-01--RF-09 & Requisitos, Diseño \\
    Descripción: estructura desacoplada y configurable & RNF-01, RNF-04 & Diseño, Construcción \\
    Descripción: prototipo con gestión de usuarios, autenticación, persistencia, CMS y flujo de pago & RF-03, RF-05, RF-06, RF-07, RF-09 & Construcción, Pruebas \\
    Alcance 1: análisis y requisitos & Todo el catálogo & Requisitos, Análisis \\
    Alcance 2: arquitectura, autenticación, roles, pagos y CMS & RF-02--RF-04, RF-06--RF-09, RNF-02 & Diseño \\
    Alcance 3: modelo de datos y mecanismos de acceso & RI-01--RI-11, RNF-02, RNF-03 & Análisis, Diseño \\
    Alcance 4: implementación del prototipo y configuración para su reutilización & RF-01--RF-13, RNF-01, RNF-05 & Construcción \\
    Alcance 5: verificación de casos de uso, integración y evaluación de la reutilización & RNF-06, RNF-01 & Pruebas \\
    Alcance 6: memoria, documentación técnica y manuales & RNF-05, RNF-07 & Despliegue, anexos \\
    \bottomrule
  \end{tabularx}
\end{table}

% ---------------------------------------------------------------------
\section{Análisis GAP}
\label{sec:req-gap}

% PENDIENTE (P-06): redacción neutra acordada con el tutor; no se nombran las
% bases de código de partida (ADR-007).
Como se justifica en el capítulo de Planificación (ADR-001), PymeKit no parte
de cero: se apoya en dos bases de código existentes, utilizadas con licencia.
La primera es una base SaaS sobre la misma pila tecnológica, que ya resolvía
la autenticación, el modelo de cuentas multi-tenant, los equipos y la
facturación. La segunda es un CMS de administración de datos para el mismo
tipo de base de datos, distribuido como aplicación de página única con su
propia API. La tabla~\ref{tab:req-gap} resume qué cubría cada una frente a
los requisitos y qué ha tenido que aportar el proyecto.

\begin{table}[htbp]
  \centering
  \footnotesize
  \caption{Brecha entre los requisitos y las bases de partida («sí»: cubierto; «parcial»: cubierto en parte; «no»: no cubierto).}
  \label{tab:req-gap}
  \begin{tabularx}{\textwidth}{@{}lccX@{}}
    \toprule
    \textbf{Req.} & \textbf{Base SaaS} & \textbf{Base CMS} & \textbf{Aportación del proyecto} \\
    \midrule
    RF-01 & parcial & no & Blog en la base de datos con lectura pública por RLS; marca e identidad propias \\
    RF-02--RF-05 & sí & no & Comprobación de sesión viva en las rutas protegidas; textos en español \\
    RF-06, RF-07 & sí & no & Integración y pruebas; configuración como módulos opcionales \\
    RF-08 & sí & no & Consola compartida con el CMS y MFA obligatorio en ella \\
    RF-09 & no & parcial & Integración total en la web, reescritura de las pantallas, consola por áreas, nombres en español, endurecimiento de seguridad \\
    RF-10, RF-11 & no & parcial & Auditoría no falsificable, autor conservado, paneles sin dependencias nuevas \\
    RF-13 & parcial & parcial & Un único sistema de traducción y toda la interfaz en español \\
    RNF-02, RNF-03 & parcial & parcial & Revisiones adversariales y corrección de brechas heredadas \\
    RNF-05 & no & no & Comentarios didácticos en español y controles automáticos del TFG \\
    \bottomrule
  \end{tabularx}
\end{table}

La brecha más importante no está en la funcionalidad, sino en la
\emph{integración} y en la \emph{seguridad}: las dos bases tenían modelos de
autorización, sistemas de formularios y de traducción y enrutadores distintos,
y el CMS de partida contenía fallos de control de acceso que se describen en
los capítulos de Diseño y Pruebas. La distinción detallada entre lo
reutilizado, lo adaptado y lo propio se presenta en el capítulo de
Construcción (sección~\ref{sec:imp-reutilizacion}).
